\section{정책금융기관으로서의 역할}

한국산업은행의 본질적 기능은 단순한 금융서비스 제공을 넘어, \textbf{정책금융(policy finance)}을 수행하는 데 있다.  
정책금융이란 시장 메커니즘만으로는 달성하기 어려운 경제적 목표를 실현하기 위해 정부 또는 공공기관이 금융을 통해 개입하는 것을 의미한다.

이 장에서는 산업은행의 역할을 시장실패 보완, 거시경제 안정화, 산업정책 수행이라는 세 가지 측면에서 분석한다.

\subsection{시장실패 보완 기능}

산업은행의 가장 기본적인 역할은 시장실패를 보완하는 것이다.

앞서 살펴본 바와 같이, 자본시장의 불완전성, 정보 비대칭성, 외부효과 등은 민간 금융만으로는 효율적인 자원 배분을 달성하기 어렵게 만든다.

이러한 상황에서 산업은행은 다음과 같은 방식으로 시장을 보완한다:

\begin{itemize}
    \item 장기·대규모 투자에 대한 자금 공급
    \item 신용위험이 높은 기업에 대한 금융 접근성 확대
    \item 사회적 수익이 높은 분야에 대한 투자 유도
\end{itemize}

이는 시장에서 과소공급되는 자원을 보충함으로써 \textbf{사회적 후생(social welfare)}을 증가시키는 효과를 가진다.

\subsection{거시경제 안정화 기능}

산업은행은 경기 변동에 대응하여 금융을 조절하는 \textbf{거시경제 안정화 장치}로 기능한다.

특히 경기 침체기에는 민간 금융기관이 위험 회피 성향을 강화하면서 대출을 축소하는 경향이 있다. 이 경우 실물경제에 대한 자금 공급이 급격히 위축되는 \textit{신용경색(credit crunch)}이 발생할 수 있다.

산업은행은 이러한 상황에서 다음과 같은 역할을 수행한다:

\begin{itemize}
    \item 기업 유동성 지원 확대
    \item 구조조정을 통한 금융 시스템 안정화
    \item 대규모 금융 공급을 통한 경기 하방 압력 완화
\end{itemize}

이러한 기능은 경기 변동을 완화하는 \textbf{완충 장치(buffer)}로서 중요한 의미를 가진다.

\subsection{산업정책 수행 기능}

산업은행은 단순한 금융기관이 아니라, 정부의 산업정책을 실행하는 핵심 수단이다.

특정 산업이나 기술 분야는 국가 경쟁력 확보를 위해 전략적으로 육성될 필요가 있으나, 민간 부문은 높은 위험과 불확실성 때문에 투자를 회피할 수 있다.

이러한 경우 산업은행은 다음과 같은 방식으로 개입한다:

\begin{itemize}
    \item 전략 산업 및 신산업에 대한 선도적 투자
    \item 산업 구조 전환을 위한 금융 지원
    \item 지역경제 및 중소기업 지원
\end{itemize}

이는 산업은행이 단순한 금융중개자가 아니라 \textbf{경제 구조를 설계하는 정책 도구}임을 의미한다.

\subsection{정책금융의 한계}

정책금융은 필요성이 존재하지만, 동시에 여러 가지 한계를 내포한다.

\subsubsection{자원 배분의 비효율성}

정부 개입은 정치적 고려나 정보 부족으로 인해 비효율적인 자원 배분으로 이어질 수 있다.  
특히 수익성이 낮은 사업에 과도한 자금이 투입될 가능성이 존재한다.

\subsubsection{도덕적 해이 (moral hazard)}

기업이 정책금융 지원을 기대하게 되면, 위험한 투자나 비효율적 경영을 지속할 유인이 발생할 수 있다.

이는 장기적으로 경제 전체의 효율성을 저해할 수 있다.

\subsubsection{민간 금융 위축 (crowding-out)}

산업은행이 시장에 적극적으로 참여할 경우, 민간 금융기관의 역할이 축소될 수 있다.  
이는 금융시장 전체의 경쟁과 혁신을 저해할 가능성을 내포한다.

\subsection{정책금융과 시장의 균형}

따라서 정책금융의 핵심 과제는 \textbf{시장과 정부 간의 적절한 균형을 유지하는 것}이다.

이 균형은 다음과 같은 원칙을 요구한다:

\begin{itemize}
    \item 시장이 실패하는 영역에 한정된 개입
    \item 민간 금융과의 보완적 관계 유지
    \item 정책 목표와 효율성 간의 균형 확보
\end{itemize}

이러한 조건이 충족될 때 정책금융은 경제 성장과 안정성에 긍정적인 기여를 할 수 있다.

\subsection{소결}

한국산업은행은 시장실패 보완, 거시경제 안정화, 산업정책 수행이라는 세 가지 핵심 기능을 통해 정책금융기관으로서 중요한 역할을 수행한다.

그러나 이러한 역할은 항상 효율성과 충돌할 가능성을 내포하고 있으며, 따라서 정책금융의 범위와 방식에 대한 지속적인 재검토가 필요하다.

결국 산업은행의 역할은 단순한 금융기관을 넘어, \textbf{국가 경제 시스템의 핵심 조정 장치}로 이해될 수 있다.